\chapter{Metodologia Geral}
\label{chap:metodologia}

\orientacao{Capítulo que o relatório intermediário não tem e o final precisa ter. Registra como o projeto foi conduzido: ritmo de trabalho, instâncias de reunião com o parceiro, forma de acompanhamento, composição das equipes por pacote e a lista de macroentregas com mês e conteúdo. Tem valor de auditoria, porque mostra a governança do projeto, e valor prático, porque a próxima equipe entende como o trabalho foi organizado.

Descreva o que de fato ocorreu. Se um rito foi abandonado no segundo ano, diga isso ou omita o ponto, mas não descreva um processo que não existiu.}

\section{Organização do Projeto}

\ph{Ritmo de trabalho, reuniões com o parceiro e forma de acompanhamento, como ocorreram.}

\section{Estrutura da Equipe}

\ph{Composição das equipes por pacote de trabalho e papel de cada uma.}

\section{Pacotes de Trabalho}

\ph{Os pacotes de trabalho, com o objetivo de cada um.}

\section{Macroentregas Executadas}

A Tabela~\ref{tab:macroentregas} lista as macroentregas do projeto.

\begin{table}[htbp]
\centering
\caption{Macroentregas executadas.}
\label{tab:macroentregas}
\estilotabela
\begin{tabularx}{\linewidth}{|C{2.9cm}|C{2.4cm}|X|}
\hline
\rowcolor{reqcolor}\cab{Macroentrega} & \cab{Mês} & \cab{Conteúdo, como no Anexo I} \\ \hline
\ph{M01} & \ph{mm/aaaa} & \ph{descrição} \\ \hline
\rowcolor{reqlight}\ph{M02} & \ph{mm/aaaa} & \ph{descrição} \\ \hline
\end{tabularx}
\end{table}
